% Partial fractions, AY2526 Final Q3(c): putting in a number, the rebuild check, the piece table, and two meanings of C.
% Copy the pattern, not the maths. \note{} is the right-hand column of an align* line.
% Equation: x^2 + x + 1 = (Ax+B)(3x+1) + C(x^2+4).

% ---- In \fast{}: the whole fast route, blocks without "Why this value?" or step titles.
\fast{\begin{gather*}
(3x+1)(x^2+4) = 3x^3 + x^2 + 12x + 4 \\
6x^3 \div 3x^3 = 2, \qquad (6x^3 + 3x^2 + 25x + 9) - 2(3x^3 + x^2 + 12x + 4) = x^2 + x + 1 \\
f(x) = 2 + \frac{x^2 + x + 1}{(3x+1)(x^2+4)}
\end{gather*}}
\ans{$f(x) = 2 + \dfrac{x^2 + x + 1}{(3x+1)(x^2+4)}$, \ for $x \neq -\dfrac13$.}
\fast{\begin{gather*}
u = \frac{x}{2}, \qquad u' = \frac12 \\
\frac{d}{dx}\tan^{-1}\!\left(\frac{x}{2}\right) = \frac{\frac12}{1 + \left(\frac{x}{2}\right)^2} = \frac{\frac12}{1 + \frac{x^2}{4}} = \frac{\frac12}{\frac{4 + x^2}{4}} = \frac{4}{2(4 + x^2)} = \frac{2}{x^2 + 4}
\end{gather*}}
\ans{$\dfrac{d}{dx}\tan^{-1}\!\left(\dfrac{x}{2}\right) = \dfrac{2}{x^2 + 4}$.}
\fast{\begin{gather*}
f(x) = 2 + \frac{x^2+x+1}{(3x+1)(x^2+4)} = 2 + \frac{Ax+B}{x^2+4} + \frac{C}{3x+1} \\
x^2 + x + 1 = (Ax+B)(3x+1) + C(x^2+4)
\end{gather*}
\par\pagebreak[0]\textbf{$x = -\frac13$.} Left side:
\begin{align*}
x^2 + x + 1 &= \tfrac19 - \tfrac13 + 1 \note{put in $x = -\tfrac13$} \\
&= \tfrac19 - \tfrac39 + \tfrac99 \note{same bottoms} \\
&= \tfrac79 \note{add the tops}
\end{align*}
\par\pagebreak[0]Right side:
\begin{align*}
(Ax+B)\cdot\chg{0} + C\left(\tfrac19 + 4\right) &= C\left(\tfrac19 + \tfrac{36}{9}\right) \note{first term is $0$} \\
&= \tfrac{37}{9}C \note{add the tops}
\end{align*}
\par\pagebreak[0]Make them equal:
\begin{align*}
\tfrac79 = \tfrac{37}{9}C \ \Rightarrow \ C &= \tfrac{7}{37} \note{divide by $\tfrac{37}{9}$}
\end{align*}
\par\pagebreak[0]\textbf{$x = 0$.}
\begin{align*}
0 + 0 + 1 &= (A\cdot 0 + B)(3\cdot 0 + 1) + C(0 + 4) \note{put in $x = 0$} \\
1 &= B + 4C \note{tidy up} \\
1 &= B + \tfrac{28}{37} \note{$C = \tfrac{7}{37}$} \\
B &= \tfrac{37}{37} - \tfrac{28}{37} = \tfrac{9}{37} \note{subtract}
\end{align*}
\par\pagebreak[0]\textbf{$x^2$-terms.}
\begin{align*}
(Ax)(3x) + Cx^2 &= (3A + C)x^2 \note{only $x^2$ terms} \\
1 &= 3A + C \note{match the $x^2$ numbers} \\
1 &= 3A + \tfrac{7}{37} \note{$C = \tfrac{7}{37}$} \\
\tfrac{30}{37} &= 3A \note{subtract $\tfrac{7}{37}$} \\
A &= \tfrac{10}{37} \note{divide by $3$}
\end{align*}
\par\pagebreak[0]\textbf{Check at $x = 2$.}
\begin{center}\renewcommand{\arraystretch}{1.6}
\begin{tabular}{lr}
 & \textbf{Value} \\ \hline
Original: $\dfrac{48 + 12 + 50 + 9}{(7)(8)} = \dfrac{119}{56}$ & $2.125$ \\ \hline
Rebuilt: $2 + \dfrac{29/37}{8} + \dfrac{7/37}{7}$ & $2.125$ \checkmark
\end{tabular}
\end{center}
\par\pagebreak[0]\textbf{Integrate each piece.} Here $x^2 + 4 = x^2 + 2^2$, so $a = 2$.
\begin{center}\small\setlength{\tabcolsep}{4pt}\renewcommand{\arraystretch}{1.8}
\begin{tabular}{lllll}
\textbf{Piece} & \textbf{Formula} & \textbf{Number in front} & \textbf{$\times$ its number} & \textbf{Result} \\ \hline
$2$ & $\int 2\,dx = 2x$ & --- & --- & $2x$ \\ \hline
$\frac{10}{37}\cdot\frac{x}{x^2+4}$ & $\frac12\ln(x^2+a^2)$ & $\frac{10}{37}$ & $\frac{10}{37}\times\frac12 = \frac{5}{37}$ & $\frac{5}{37}\ln(x^2+4)$ \\ \hline
$\frac{9}{37}\cdot\frac{1}{x^2+4}$ & $\frac1a\tan^{-1}\frac{x}{a}$ & $\frac{9}{37}$ & $\frac{9}{37}\times\frac12 = \frac{9}{74}$ & $\frac{9}{74}\tan^{-1}\frac{x}{2}$ \\ \hline
$\frac{7}{37}\cdot\frac{1}{3x+1}$ & $\frac1a\ln|ax+b|$, $a = 3$ & $\frac{7}{37}$ & $\frac{7}{37}\times\frac13 = \frac{7}{111}$ & $\frac{7}{111}\ln|3x+1|$
\end{tabular}
\end{center}
\par\pagebreak[0]\textbf{Add the results.}
\[ \int f\,dx = 2x + \tfrac{5}{37}\ln(x^2+4) + \tfrac{9}{74}\tan^{-1}\!\tfrac{x}{2} + \tfrac{7}{111}\ln|3x+1| + C \]
This $+C$ is the constant of integration. It is not the $C = \frac{7}{37}$.}

% ---- In \whyworks: the same blocks inside steps, each step opening with its reason.
\step{\choice Find $B$ by putting in $x = 0$.}
\reason{Why this value? Every term with an $x$ in it becomes $0$, so $A$ drops out.}
\begin{align*}
0 + 0 + 1 &= (A\cdot 0 + B)(3\cdot 0 + 1) + C(0 + 4) \note{put in $x = 0$} \\
1 &= B + 4C \note{tidy up} \\
1 &= B + \tfrac{28}{37} \note{$C = \tfrac{7}{37}$} \\
B &= \tfrac{37}{37} - \tfrac{28}{37} = \tfrac{9}{37} \note{subtract}
\end{align*}

% ---- Rebuild check: Original and Rebuilt at a test number, value column, tick.
\step{Check the rebuilt $f$ on a number: put $x = 2$.}
\reason{If the split is right, the original $f$ and the rebuilt $2 + R(x)$ give the same value at any $x$.}
\begin{center}\renewcommand{\arraystretch}{1.6}
\begin{tabular}{lr}
 & \textbf{Value} \\ \hline
Original: $\dfrac{48 + 12 + 50 + 9}{(7)(8)} = \dfrac{119}{56}$ & $2.125$ \\ \hline
Rebuilt: $2 + \dfrac{29/37}{8} + \dfrac{7/37}{7}$ & $2.125$ \checkmark
\end{tabular}
\end{center}

% ---- Piece table: each piece's number in front times its formula's number.
\step{Collect the four pieces in one table.}
\reason{Each piece is a number in front times a formula. The formula brings its own number ($\frac12$ or $\frac13$). Multiply the two numbers to get the piece's final number. Here $x^2 + 4 = x^2 + 2^2$, so $a = 2$.}
\begin{center}\small\setlength{\tabcolsep}{4pt}\renewcommand{\arraystretch}{1.8}
\begin{tabular}{lllll}
\textbf{Piece} & \textbf{Formula} & \textbf{Number in front} & \textbf{$\times$ its number} & \textbf{Result} \\ \hline
$2$ & $\int 2\,dx = 2x$ & --- & --- & $2x$ \\ \hline
$\frac{10}{37}\cdot\frac{x}{x^2+4}$ & $\frac12\ln(x^2+a^2)$ & $\frac{10}{37}$ & $\frac{10}{37}\times\frac12 = \frac{5}{37}$ & $\frac{5}{37}\ln(x^2+4)$ \\ \hline
$\frac{9}{37}\cdot\frac{1}{x^2+4}$ & $\frac1a\tan^{-1}\frac{x}{a}$ & $\frac{9}{37}$ & $\frac{9}{37}\times\frac12 = \frac{9}{74}$ & $\frac{9}{74}\tan^{-1}\frac{x}{2}$ \\ \hline
$\frac{7}{37}\cdot\frac{1}{3x+1}$ & $\frac1a\ln|ax+b|$, $a = 3$ & $\frac{7}{37}$ & $\frac{7}{37}\times\frac13 = \frac{7}{111}$ & $\frac{7}{111}\ln|3x+1|$
\end{tabular}
\end{center}

% ---- One letter, one meaning: say which C is which.
\step{Add the results, and add $+C$.}
\reason{This $+C$ is the constant of integration. It is not the $C = \frac{7}{37}$ from the partial fractions.}
\[ \int f(x)\,dx = 2x + \frac{5}{37}\ln(x^2+4) + \frac{9}{74}\tan^{-1}\!\left(\frac{x}{2}\right) + \frac{7}{111}\ln|3x+1| + C \]
